%! Regression Lines and Residuals — Unit 2, Lesson 2.3 — Student Packet Cover Sheet
\documentclass[12pt]{article}
\usepackage{apstats-article}
\usepackage{apstats-boxes}
\usepackage{ltablex}
\keepXColumns

\begin{document}

% Full-bleed navy banner, sized to the title block by \coverbanner.
\coverbanner{Unit 2}{Lesson 2.3 \quad Regression Lines and Residuals}

\namedateperiod
\vspace{0.18in}

\boxguard[14]
\begin{learningtargetbox}
\small
% GRADUAL RELEASE: the vocabulary is taught in the guided notes, so the targets name it
% plainly. The AP Learning Objectives (DAT-1.D, DAT-1.E, DAT-1.F) stay in the plan.
\begin{enumerate}[leftmargin=1.5em, itemsep=5pt]
  \item \ldots{} use a \textbf{regression line} $\hat{y} = a + bx$ to find a
        \textbf{predicted value} $\hat{y}$ in context, and tell it apart from the actual
        value $y$.
  \item \ldots{} recognize \textbf{extrapolation} and explain why a prediction far outside the
        data is not trustworthy.
  \item \ldots{} calculate a \textbf{residual} as $y - \hat{y}$ and say in context whether the
        line predicted too high or too low.
  \item \ldots{} read a \textbf{residual plot} and decide whether a linear model is
        appropriate.
\end{enumerate}
\end{learningtargetbox}

\vspace{0.14in}

\boxguard[14]
\begin{tocbox}
% A CONTENTS LIST ONLY: no Score column and no Total row. Whether a component counts is said in
% its description.
\small
\renewcommand{\arraystretch}{1.5}
\begin{tabularx}{\linewidth}{@{} c l X @{}}
\rowcolor{sky}
\textbf{\#} & \textbf{Component} & \textbf{Description} \\
\midrule
1 & Warm-Up & Plugging into a line, and a point above or below it \\
2 & Guided Notes \& Practice & Predicting pizza-delivery times, how far to trust them, and
                 what the misses tell us \\
3 & AP Practice & A coffee shop's hot-drink sales: two multiple-choice items and one
                 free-response set --- practice, not a grade \\
4 & Homework & Used-car prices: predictions, residuals, and a residual plot ---
                 \textbf{scored, due the first class after two study halls} \\
\end{tabularx}
\end{tocbox}

\vspace{0.14in}

\boxguard[8]
\begin{remindbox}
\small
A regression line gives a \textbf{prediction}, $\hat{y}$; the data give what actually
happened, $y$. They rarely agree, and the gap is the \textbf{residual} --- always
\textbf{actual minus predicted}, in that order. A positive residual means the real value came
out \emph{bigger} than predicted: the point sits above the line, and the line guessed too
\emph{low}. And a line only speaks for the $x$-values it was built from --- step outside them
and the prediction is a guess.
\end{remindbox}

\end{document}
